\documentclass[11pt]{article}
\usepackage[margin=1in]{geometry}
\usepackage{amsmath,amssymb,amsthm,mathtools}
\usepackage{graphicx}
\usepackage{float}
\usepackage{microtype}
\usepackage[hidelinks]{hyperref}

\newtheorem{theorem}{Theorem}
\newtheorem{corollary}[theorem]{Corollary}
\newtheorem{proposition}[theorem]{Proposition}
\newcommand{\C}{\mathbb C}
\newcommand{\AT}{A_T}

\title{Two Derivation Obstructions for Compact\\Quasinilpotent Generators}
\author{Partial-result packet for arXiv:1008.5241}
\date{12 August 2026}

\begin{document}
\maketitle

\begin{center}
\fbox{\begin{minipage}{0.9\textwidth}
\textbf{Status: candidate substantial partial result, likely valid; human
review requested.}
For \(\AT=\overline{\operatorname{span}}\{T^n:n\ge1\}\), weak amenability
forces both \(\AT=\overline{\AT^2}\) and
\([B(H),T]\cap\AT=\{0\}\).  These necessary conditions rule out every
nonzero compact unilateral weighted shift and every fractional-integration
operator \(J_{1/r}\), including the classical Volterra operator.  They do
not decide whether an exceptional compact quasinilpotent satisfying both
conditions can be weakly amenable.
\end{minipage}}
\end{center}

\section{The source question}

For \(T\in B(H)\), put
\[
 \AT=\overline{\operatorname{span}}\{T^n:n\ge1\}.
\]
This is a commutative, possibly nonunital Banach algebra.  The source calls
\(T\) weakly amenable when \(\AT\) is weakly amenable.  Remark~2.3 asks:
\begin{quote}
Can a compact quasinilpotent operator be weakly amenable?
\end{quote}
The same remark's question about compact normal operators is answered later
in the source, but the quasinilpotent question is left open \cite{SWJ}.

\begin{figure}[H]
\centering
\includegraphics[width=0.96\textwidth]{figures/open_problem_crop.png}
\caption{Source Remark~2.3 and its Volterra discussion (PDF page 4).}
\end{figure}

For commutative Banach algebras we use the standard characterization quoted
by the source: weak amenability is equivalent to the vanishing of every
bounded derivation into every symmetric Banach bimodule.  This formulation
makes two elementary but effective obstructions visible.

\section{The two obstructions}

\begin{theorem}[Power-span and commutator obstructions]
\label{thm:obstructions}
Let \(0\ne T\in B(H)\).  If either
\begin{enumerate}
\item \(T\notin\overline{\operatorname{span}}\{T^n:n\ge2\}\), or
\item there is \(X\in B(H)\) such that
      \(0\ne[X,T]\in\AT\),
\end{enumerate}
then \(\AT\) is not weakly amenable.
Consequently, every possible positive answer to the source question must
satisfy
\[
 \AT=\overline{\AT^2},
 \qquad [B(H),T]\cap\AT=\{0\}.
 \tag{1}
\]
\end{theorem}

\begin{proof}
First observe that
\[
 \overline{\AT^2}
 =\overline{\operatorname{span}}\{T^n:n\ge2\}.
 \tag{2}
\]
Under condition 1, Hahn--Banach gives a bounded functional
\(\delta:\AT\to\C\) such that \(\delta(T)=1\) and
\(\delta(\overline{\AT^2})=0\).  Make \(\C\) a symmetric \(\AT\)-bimodule
with the trivial actions.  Then for \(a,b\in\AT\),
\[
 \delta(ab)=0=a\cdot\delta(b)+\delta(a)\cdot b.
\]
Thus \(\delta\) is a nonzero bounded derivation into a symmetric module, so
\(\AT\) is not weakly amenable.

Under condition 2, write \(q=[X,T]\).  Since \(q\in\AT\), it commutes with
\(T\).  Hence, for every polynomial \(p\) with \(p(0)=0\),
\[
 [X,p(T)]=p'(T)q\in\AT. \tag{3}
\]
The commutator map is bounded on \(B(H)\), so density and the closedness of
\(\AT\) extend it to a bounded derivation
\(D:\AT\to\AT\), \(D(a)=[X,a]\).  The target \(\AT\), with its usual
multiplication actions, is symmetric; moreover \(D(T)=q\ne0\).  Again weak
amenability fails.  Finally, (2) turns the negation of condition 1 into the
first assertion of (1), and the negation of condition 2 into the second.
\end{proof}

\paragraph{Proof intuition.}
Weak amenability forbids every continuous first-order differential on
\(\AT\).  If the first power survives modulo products, a scalar functional
detecting it is already such a differential.  If an ambient commutator
lands back in the power algebra, the usual operator commutator restricts to
a differential on \(\AT\).  A positive example must evade both mechanisms.

\section{Homogeneous generators and weighted shifts}

\begin{proposition}[Gauge obstruction]
\label{prop:gauge}
Suppose \(\AT\) admits a strongly continuous uniformly bounded circle action
\((\gamma_\zeta)_{|\zeta|=1}\) satisfying
\(\gamma_\zeta(T)=\zeta T\).  Then \(\AT\) is not weakly amenable.
\end{proposition}

\begin{proof}
The first Fourier projection
\[
 P_1(a)=\int_{\mathbb T}\zeta^{-1}\gamma_\zeta(a)\,dm(\zeta)
\]
is bounded.  It satisfies \(P_1(T)=T\) and \(P_1(T^n)=0\) for \(n\ge2\).
Therefore \(T\) cannot belong to the closed span of the higher powers, and
Theorem~\ref{thm:obstructions}(1) applies.
\end{proof}

\begin{corollary}[Compact unilateral weighted shifts]
Every nonzero compact unilateral weighted shift is quasinilpotent and is not
weakly amenable.
\end{corollary}

\begin{proof}
Let \(We_n=w_ne_{n+1}\) on \(\ell^2(\mathbb N_0)\).  Compactness is
equivalent to \(w_n\to0\).  Given \(\varepsilon>0\), all but finitely many
factors in every long consecutive product of weights have modulus at most
\(\varepsilon\); hence \(\limsup_k\|W^k\|^{1/k}\le\varepsilon\), and
\(W\) is quasinilpotent.

For \(|\zeta|=1\), define \(U_\zeta e_n=\zeta^ne_n\).  Then
\(U_\zeta WU_\zeta^*=\zeta W\).  Conjugation gives the circle action in
Proposition~\ref{prop:gauge}; norm continuity on \(\AT\) follows first for
polynomials and then by uniform approximation.  Thus \(W\) is not weakly
amenable.
\end{proof}

\section{Fractional integration, including Volterra}

For \(\alpha>0\), let \(J_\alpha\) be the Riemann--Liouville operator on
\(L^2(0,1)\):
\[
 (J_\alpha f)(x)=\frac1{\Gamma(\alpha)}
 \int_0^x(x-t)^{\alpha-1}f(t)\,dt. \tag{4}
\]
The case \(\alpha=1\) is the classical Volterra operator.

\begin{theorem}[A fractional Volterra family]
\label{thm:RL}
For every integer \(r\ge1\), the compact quasinilpotent operator
\(J_{1/r}\) is not weakly amenable.
\end{theorem}

\begin{proof}
We first record compactness and quasinilpotence.  Truncate the convolution
kernel \(s^{\alpha-1}/\Gamma(\alpha)\) to \(s\ge\varepsilon\).  The
resulting kernel on \((0,1)^2\) is square integrable, so its operator is
Hilbert--Schmidt.  Young's inequality bounds the discarded operator by
\[
 \frac1{\Gamma(\alpha)}\int_0^\varepsilon s^{\alpha-1}\,ds
 =\frac{\varepsilon^\alpha}{\Gamma(\alpha+1)}.
\]
Thus \(J_\alpha\) is a norm limit of compact operators.

The beta-integral identity gives the semigroup law
\(J_\alpha J_\beta=J_{\alpha+\beta}\).  Young's inequality also gives
\[
 \|J_\alpha^n\|=\|J_{n\alpha}\|
 \le \frac1{\Gamma(n\alpha+1)}.
\]
Stirling's formula now yields
\(\lim_n\|J_\alpha^n\|^{1/n}=0\).

Let \(M_x\) denote multiplication by \(x\).  A direct kernel calculation
from (4) gives
\[
 [M_x,J_\alpha]=\alpha J_{\alpha+1}. \tag{5}
\]
If \(\alpha=1/r\), then the semigroup law turns (5) into
\[
 [M_x,J_{1/r}]=\frac1r J_{1/r}^{\,r+1}
 \in A_{J_{1/r}}\setminus\{0\}.
\]
Theorem~\ref{thm:obstructions}(2) completes the proof.
\end{proof}

In particular, \(J_1\) is not weakly amenable.  The source records the
Volterra operator as character amenable; the conclusion here is consistent
with that statement and shows that character amenability is strictly weaker
for this canonical compact quasinilpotent.

\section{Boundary of the result and upgrade audit}

Theorem~\ref{thm:obstructions} gives necessary, not sufficient, conditions.
A full negative solution would follow from a universal theorem that every
nonzero compact quasinilpotent either fails
\(\AT=\overline{\AT^2}\) or admits a nonzero commutator
\([X,T]\in\AT\).  Neither step is presently justified.  Volterra-type
power algebras show that the first alternative need not be automatic, while
the second asks for exact membership in the range of the Rosenblum operator
\(X\mapsto[X,T]\), not merely approximate triangular information.

Eight focused upgrade attempts covered the current literature, nilpotent and
finite-type reductions, transfer of the abstract radical weakly amenable
example, normalized powers, the essentiality obstruction, gauge actions,
commutator derivations, and the universal Rosenblum/triangularization route.
The last two successful mechanisms produce the two broad classes above; the
remaining universal step has no credible closure argument.  The full source
question therefore remains open in this packet.

\paragraph{Verification and novelty.}
The proof is self-contained apart from the standard symmetric-module
characterization of weak amenability quoted by the source.  Human review
should focus on the range assertion (3), the norm-continuity of the gauge
action, and the kernel calculation (5).  On 12 August 2026, all cheap run
indexes and exact/current literature searches were checked.  Later work on
singly generated radical operator algebras treats compact weighted shifts
and Volterra-type algebras but does not decide weak amenability \cite{Peters}.
An earlier abstract construction gives a singly generated radical weakly
amenable Banach algebra, but on a Banach space isometric to \(\ell^1\), not
as the Hilbert-space compact-operator example asked for here \cite{LRRW}.
Novelty confidence is moderate and provisional.

\paragraph{Human-review recommendation.}
Review as a substantial partial result.  It materially narrows any possible
positive example and settles two natural, infinite families, but it is not a
proof or counterexample to the existence question.

\begin{thebibliography}{9}

\bibitem{SWJ}
L.~Y.~Shi, Y.~J.~Wu, and Y.~Q.~Ji,
\emph{An example of weakly amenable and character amenable operator},
arXiv:1008.5241; Illinois J. Math. 55 (2011), no.~4, 1415--1422.

\bibitem{LRRW}
R.~J.~Loy, C.~J.~Read, V.~Runde, and G.~A.~Willis,
\emph{Amenable and weakly amenable Banach algebras with compact
multiplication}, J. Funct. Anal. 171 (2000), 78--114.

\bibitem{Peters}
J.~Peters,
\emph{Singly generated radical operator algebras},
arXiv:2309.13530; New York J. Math. 30 (2024), 398--421.

\end{thebibliography}

\end{document}
